\section{A Second Domain: Typed Skill Routing}\label{sec:routing}
Every study above measures generation recall on one corpus. Whether the instrument transfers is a separate question, so we applied it to a task of a different shape: a discriminative choice over a curated set, where the model is shown every candidate and must name one, or decline.

\paragraph{Task and exposure.} An agent runtime routes each user turn to one of 115 \emph{skills}, each described by a short rubric stating what it is for and when it must not be chosen. The router puts one typed choice to a judge: the turn as state, the 115 rubrics as options, plus a \texttt{none} option for turns no skill should handle. The gold label for a turn is the key of one exposed rubric, or \texttt{none}. Exposure is therefore total: the gold answer is, by construction, among the text the judge is shown. This is the condition under which the paper argues a copy ceiling is mandatory, and it is stronger here than in Study~1, where $3.6\%$ of gold lay outside the scaffold.

\paragraph{Adapting the ceiling to a choice task.} A no-op extractor has no meaning when the answer is a key rather than a passage, so the copy procedure becomes the best \emph{judge-free surface match} over the exposed options: rank the 115 rubrics against the turn and take the top. We use two rankers with no trained judgement in them, BM25 over the rubric text ($k_1{=}1.5$, $b{=}0.75$, a fixed stoplist) and cosine similarity under a 384-dimensional sentence embedding. The naive ceiling concedes \texttt{none} to the judge, since a ranker has no rubric to match a decline against; that concession supplies a large and spurious gain, so the reported ceiling is the \emph{fair} one, in which each ranker is also given a decline rule and its threshold is tuned to maximise the ranker's own accuracy on the evaluated corpus. The ceiling is thereby flattered, the judge's gain is a floor, and the ranker is held to the same affordance the judge has.

\paragraph{Judge.} The judge is a 4B-parameter natural-language-inference cross-encoder~\cite{openjev2026}\footnote{\texttt{AlexWortega/openjev}, checkpoint \texttt{qwen3.5-4b-nli-v2}, revision pinned; MIT.} scoring each (turn, rubric) pair independently as entailment, with \texttt{none} answered when the best score falls below a threshold of $0.5$; a threshold sweep from $0.05$ to $0.7$ moves top-1 accuracy by $5.8$ points with a flat optimum between $0.25$ and $0.40$, so the result is not balanced on that constant. It replaced a 421M-parameter encoder whose options share a single 256-token head budget and are each truncated to 48 tokens: that engine scored $53.5\%$ top-1 and its decline behaviour collapsed as options were added, the failure that motivated independent scoring and uncompressed rubrics. We wrote the corpus: 86 labelled turns (16 \texttt{none}, 70 skill), with each turn tagged by the class of confusion it was written to probe and by whether the rubric's discriminating clause sits late in its text, beyond the 48-token cut. The judge is not called directly: it sits behind a typed-decision façade\footnote{Sovereign System One, in the agent runtime \url{https://github.com/DreamLab-AI/agentbox}: a Jev-compatible wire protocol (choice, score, probability) with a local engine behind it; the routing corpus, the eval rig and both engines are released there.} that speaks the same wire protocol as the cloud judge it replaces, so the router's request was byte-identical across cloud, the superseded encoder and the cross-encoder, and the engine is a URL chosen by this measurement. Both the corpus and the rig are released with the runtime.

\paragraph{Result.} Table~\ref{tab:routing} and Figure~\ref{fig:routing} give the outcome. Against the naive ceilings (BM25 $62.8\%$ top-1, cosine $62.8\%$, neither able to decline) the judge's gain is over $+25$ points; against the fair lexical ceiling it is $+11.6$ top-1 and $+12.8$ top-3, with a signed per-item mean of $+0.116$ and a $95\%$ interval of $[+0.012,+0.221]$: sixteen turns the judge alone gets right against six the ranker alone gets right. The interval excludes zero, but only just, and it holds from about $n=70$ against the 86 turns we have, so this is a narrow margin rather than a comfortable one. Against the fair embedding ceiling the gain is $+20.9$ top-1 with an interval of $[+0.106,+0.312]$. The lexical ranker, which needs no model, remains the strong baseline: $76.7\%$ top-1 and $86.0\%$ top-3 at negligible cost, against the judge's $88.4\%$ and $98.8\%$ at about five seconds per turn.

\begin{table*}[t]\centering\small
\caption{Skill routing, $n=86$, top-1 accuracy by population. The ceiling columns give a judge-free
ranker's accuracy and the judge's gain over it; both rankers carry an oracle-tuned decline rule
(BM25 threshold $9.16$, fired on 17 turns; cosine threshold $0.59$, fired on 6). ``Late'' marks turns
whose gold rubric's discriminating clause lies beyond a 48-token cut.}
\label{tab:routing}
\begin{tabular}{lrrrrrr}
\toprule
Population & $n$ & Judge & BM25 & Gain & Cosine & Gain \\
\midrule
All turns              & 86 & 88.4 & 76.7 & \gainpos{11.6} & 67.4 & \gainpos{20.9} \\
Skill turns            & 70 & 85.7 & 77.1 & \gainpos{8.6}  & 75.7 & \gainpos{10.0} \\
\texttt{none} turns    & 16 & 100.0 & 75.0 & \gainpos{25.0} & 31.2 & \gainpos{68.8} \\
Late discriminator     & 32 & 96.9 & 78.1 & \gainpos{18.8} & 93.8 & \gainpos{3.1}  \\
Early discriminator    & 54 & 83.3 & 75.9 & \gainpos{7.4}  & 51.9 & \gainpos{31.5} \\
Boundary class         & 12 & 83.3 & 50.0 & \gainpos{33.3} & 58.3 & \gainpos{25.0} \\
Near-neighbour class   & 50 & 86.0 & 82.0 & \gainpos{4.0}  & 76.0 & \gainpos{10.0} \\
\bottomrule
\end{tabular}
\end{table*}

\begin{figure*}[t]\centering
\begin{tikzpicture}
\begin{axis}[xbar,width=13cm,height=6.9cm,bar width=9pt,xmin=-4,xmax=38,
  xlabel={Judge gain over the fair lexical copy ceiling (percentage points, top-1)},
  enlarge y limits=0.08, ytick=data,
  symbolic y coords={early discriminator,near-neighbour,skill items,all items,boundary,late discriminator,single,none items},
  yticklabel style={font=\small}, xmajorgrids, grid style={grayln},
  nodes near coords, nodes near coords style={font=\footnotesize,/pgf/number format/showpos},
  nodes near coords align={horizontal},
  axis line style={ink!60}]
\addplot[fill=teal0!25,draw=teal0] coordinates {
  (7.4,early discriminator) (4.0,near-neighbour) (8.6,skill items) (11.6,all items)
  (33.3,boundary) (18.8,late discriminator) (20.0,single) (25.0,none items)};
\draw[burnt,dashed,line width=0.7pt] (axis cs:0,early discriminator) -- (axis cs:0,none items);
\end{axis}\end{tikzpicture}
\caption{Typed skill routing, $n=86$: where the judge's gain over the fair copy ceiling lies.
Each bar is judge top-1 accuracy minus the accuracy of a BM25 ranker over the same 115 exposed
rubrics, with that ranker granted a decline rule and its threshold tuned on this very corpus, so the
ceiling is flattered and every bar is a floor. The gain is not spread: it concentrates on declining
(\texttt{none}, $+25.0$), on rubrics whose discriminating clause sits late ($+18.8$) and on boundary
cases ($+33.3$), each a consequence of a stated architectural choice (independent scoring; uncompressed
rubrics), while ordinary near-neighbour picking among exposed options gains $+4.0$. The all-items interval is $+0.116$,
$[+0.012,+0.221]$, clear of zero by one turn's margin; against an embedding ceiling it is $+0.209$,
$[+0.106,+0.312]$.}
\label{fig:routing}
\end{figure*}


\paragraph{Threshold, escalation and confidence in detail.} Three per-turn artefacts of the run are worth tabulating because each settles a question the aggregate cannot. Table~\ref{tab:routing-sweep} sweeps the judge's decline threshold over one full corpus pass per value: top-1 accuracy moves $5.8$ points across $0.05$--$0.7$, \texttt{none} recall is $100\%$ from $0.25$ to $0.4$ and falls away below $0.25$, and precision erodes above $0.4$, so the optimum is a plateau from $0.25$ to $0.40$ with $0.3$ at its centre. The deployed threshold of $0.5$, which every figure in this section uses, sits just outside that plateau and costs $1.1$ points: the reported result is not the best this judge can do on this corpus, which is the conservative direction for a threshold we did not tune before measuring. Table~\ref{tab:routing-cascade} escalates to the judge only when the lexical ranker's top-two BM25 margin is below a cut: at $40\%$ escalation the cascade lands within one turn of the judge's accuracy for two-fifths of its mean latency, and the curve is not monotone at this $n$, which is itself a caution against reading any single row as a tuned operating point. The ranker and judge disagree on $25.6\%$ of turns and their union is right on $95.3\%$, so the complementarity is real whatever the cascade is worth. Table~\ref{tab:routing-calib} bins the judge's reported confidence: accuracy rises monotonically with confidence, the engine is right $83\%$ of the time when claiming $0.34$, no wrong answer carries confidence above $0.548$, and the rank separation of right from wrong is $0.728$ (AUC). The judge is under-confident but rank-informative, which is the opposite of the over-confident failure that makes confidence gating unsafe in general, and which we had measured on a different judge for this same task: there, a wrong pick carried $0.94$. Calibration is evidently a property of the judge rather than of the task, so a gating rule fitted to one does not transfer unexamined to another; six turns above $0.7$ make this a direction to test rather than a licence to gate.

\begin{table*}[t]\centering\small
\caption{Judge decline threshold, one full corpus pass per row ($n=86$). The optimum is a plateau from
$0.25$ to $0.40$ at $89.5\%$; the deployed value ($0.5$) sits just past it, costing $1.1$ points.}
\label{tab:routing-sweep}
\begin{tabular}{rrrrrr}
\toprule
Threshold & Top-1 & Top-3 & \texttt{none} recall & \texttt{none} precision & p50 ms \\
\midrule
0.05 & 83.7 & 100.0 & 62.5  & 100.0 & 4998 \\
0.10 & 88.4 & 98.8  & 87.5  & 93.3  & 4996 \\
0.20 & 88.4 & 98.8  & 93.8  & 88.2  & 4998 \\
0.25 & 89.5 & 98.8  & 100.0 & 88.9  & 4996 \\
0.30 & 89.5 & 98.8  & 100.0 & 88.9  & 4995 \\
0.40 & 89.5 & 98.8  & 100.0 & 84.2  & 4991 \\
0.50 & 88.4 & 98.8  & 100.0 & 80.0  & 4995 \\
0.60 & 88.4 & 98.8  & 100.0 & 80.0  & 4981 \\
0.70 & 87.2 & 97.7  & 100.0 & 76.2  & 4994 \\
\bottomrule
\end{tabular}
\end{table*}

\begin{table}[t]\centering\small
\caption{Escalation cascade: answer with the BM25 ranker unless its top-two margin is below the cut,
then call the judge. Mean latency counts only judged turns (the ranker's cost is negligible).}
\label{tab:routing-cascade}
\begin{tabular}{rrrrr}
\toprule
Margin cut & Escalated & Top-1 & Mean ms & vs judge \\
\midrule
0.02  & 0\%   & 76.7 & 0    & $-11.6$ \\
0.97  & 20\%  & 80.2 & 954  & $-8.1$  \\
1.51  & 28\%  & 84.9 & 1319 & $-3.5$  \\
2.87  & 40\%  & 87.2 & 1896 & $-1.2$  \\
4.72  & 50\%  & 86.0 & 2395 & $-2.3$  \\
13.72 & 78\%  & 89.5 & 3738 & $+1.2$  \\
--    & 100\% & 88.4 & 4830 & $0.0$   \\
\bottomrule
\end{tabular}
\end{table}

\begin{table}[t]\centering\small
\caption{Judge confidence against outcome. Mean confidence is $0.457$ on correct answers and $0.323$
on wrong ones; the highest confidence on any wrong answer is $0.548$.}
\label{tab:routing-calib}
\begin{tabular}{lrrrr}
\toprule
Confidence bin & $n$ & Accuracy & Mean conf. & Gap \\
\midrule
$[0.00,0.50)$ & 53 & 83.0  & 0.341 & $+0.489$ \\
$[0.50,0.70)$ & 27 & 96.3  & 0.564 & $+0.399$ \\
$[0.70,0.85)$ & 5  & 100.0 & 0.770 & $+0.230$ \\
$[0.85,0.95)$ & 1  & 100.0 & 0.856 & $+0.144$ \\
\bottomrule
\end{tabular}
\end{table}

\paragraph{Where the gain lies.} The decomposition is the finding. Over the fair lexical ceiling, the judge's advantage concentrates on declining ($+25.0$ on \texttt{none} turns), on late-discriminator turns ($+18.8$) and on boundary cases ($+16.7$), and is small on near-neighbour picking ($+4.0$) and early-discriminator turns ($+3.7$). Each of the three large gains is the predicted consequence of a design choice rather than of general model quality: independent scoring lets \texttt{none} be a threshold instead of a competitor for mass, and uncompressed rubrics preserve the ``not for'' clause that decides late and boundary cases. On ordinary matching among exposed options, the ranker and the judge are close, which is what total exposure predicts. Read this way, the ceiling in this domain identifies what the judge adds and shows the remainder to be delivery of exposed text.

\paragraph{Disagreement as a label audit.} The two procedures err on different turns, and their union is right on $95.3\%$ against $88.4\%$ for the judge alone. Before re-adjudication (below) seven turns were ones the ranker got right and the judge did not; all seven are near-miss sibling pairs, and one was a corpus error rather than a judge error: for \emph{``Take a screenshot of the staging login page and fill in the demo credentials''} the corpus says \texttt{browser}, the judge says \texttt{browser-automation}, and the \texttt{browser} rubric itself reads ``for a browser task choose \texttt{browser-automation}, not this''. The judge followed the rubric's stated boundary; the ranker won on vocabulary overlap. A ranker wins by surface, a judge by the rubric's semantics, so systematic disagreement between them localises either a mislabelled item or a genuinely ambiguous one. When gold derives from exposed text, the copy control is therefore also a cheap pointer at the rows to re-adjudicate, a use the ceiling's definition does not anticipate.

\paragraph{Re-adjudication, and what it did to the headline.} The disagreement audit was run before the headline was fixed, and it changed it, so the sequence and the effect are both reported here. Re-reading all seven disputed turns against the rubrics' own stated boundaries, we corrected one label: turn 14, from \texttt{browser} to \texttt{browser-automation}, on the criterion quoted above, which is a clause in the rubric rather than a property of any model's answer. The other six we left standing against the judge; four are genuinely ambiguous (a request naming both synthesis and authoritative government sources matches two rubrics almost verbatim) and two are plain judge errors, where it declined against a rubric that matches the turn word for word.

That single correction moved the headline from a margin that crossed zero to one that does not: signed gain against the fair lexical ceiling went from $+0.093$, $[-0.013,+0.199]$, to $+0.116$, $[+0.012,+0.221]$, and judge accuracy from $87.2\%$ to $88.4\%$ while the lexical ceiling fell from $77.9\%$ to $76.7\%$. A single item deciding significance is the precise hazard of adjudicating labels on a corpus this small, and a reader should treat the pre-adjudication figures as the conservative ones. Three things limit the damage and none of them removes it: the adjudication criterion was the rubric's text, which is independent of which system answered; the audit corrected one of seven and left six against the judge, including the two it plainly got wrong; and the direction of the error was predicted by the mechanism before it was measured, since a ranker wins such a turn on vocabulary overlap precisely when the rubric's routing clause says it should not. The durable claim here is not the corrected interval. It is that a judge-free control surfaced a labelling error in a corpus its authors had already reviewed, and that on a corpus of this size the same audit is worth more than the significance it happens to produce.

\paragraph{What this does and does not show.} The instrument transferred to a discriminative task with a different judge architecture and returned a signed, localised answer rather than a uniform deflation. Its magnitudes do not transfer: the corpus is small, single-seed and written by us, who also wrote the rubrics it is scored against, so its phrasing echoes them and flatters a lexical ranker; the headline interval against that ranker is underpowered; and the \texttt{none} concession, fair-ceiling threshold and late-discriminator tag are all choices we made and report rather than standards. What the study licenses is a claim about method: a choice task whose gold is among its shown options should report a fair copy ceiling beside accuracy, and read the gain by population rather than in aggregate.

\paragraph{Scope of the engine-side tables.} Tables~\ref{tab:routing-sweep}--\ref{tab:routing-calib} concern the judge and its deployment rather than the copy ceiling, and are reported because the per-turn control produced them at no extra cost. Neither the confidence result nor the cascade bears on the instrument; both rest on the same 86 turns.

